\chapter*{Seminario de investigación: las preguntas comienzan hoy}
\addcontentsline{toc}{chapter}{Seminario de investigación y proyecto}
\begin{investigacion}[Meta intelectual de la última hora]
No se pretende seleccionar un tema definitivo ni entregar un experimento acabado.
Se pretende aprender a distinguir una demostración tecnológica de una pregunta
contrastable y a identificar \textbf{dos áreas de interés}, con un párrafo razonado
por área. La ficha formal será en S05; la línea base medida, en S06; y el
prerregistro, en S07. Esta pauta procede del documento de investigación del curso,
\emph{Integración e investigación aplicada}, capítulo 3 (2026).
\end{investigacion}

\section*{Problema inicial: tokenización de lenguaje técnico en español}
\addcontentsline{toc}{section}{Una pregunta científica nacida de la teoría}
Una herramienta puede fragmentar \emph{empacadura}, \emph{revestidor},
\emph{pérdida de circulación} o \emph{9 5/8 pulgadas} de maneras distintas.
Esta observación sugiere una pregunta, pero aún no constituye una hipótesis
científica: ¿una tokenización especializada necesita menos unidades por término,
y qué costo tiene ese ahorro en una predicción probabilística razonable?

La matriz institucional identifica este hilo como \textbf{IC10}: penalización de
la tokenización en textos técnicos en español. Los observables propuestos son
\emph{tokens por palabra} y \emph{bits por carácter}; el catálogo vincula la
primera exploración con D01 y D02. El documento proporciona \emph{referencias de
fuentes de datos}, no un corpus descargado y preparado para la clase.

\begin{figure}[H]\centering
\begin{tikzpicture}[>=Stealth,font=\small,
 q/.style={draw=azul,fill=azulclaro,rounded corners,align=center,text width=3.2cm,minimum height=13mm}]
\node[q] (p) at (0,0) {\textbf{Pregunta}\\¿Qué se fragmenta?};
\node[q] (h) at (4.2,0) {\textbf{Hipótesis}\\un dominio cambia la tasa};
\node[q] (m) at (8.4,0) {\textbf{Medida}\\tokens/palabra; BPC};
\node[q] (d) at (8.4,-2.0) {\textbf{Datos}\\corpus con licencia};
\node[q] (v) at (4.2,-2.0) {\textbf{Validación}\\documentos reservados};
\node[q] (r) at (0,-2.0) {\textbf{Resultado}\\efecto e incertidumbre};
\draw[->,thick] (p)--(h);\draw[->,thick] (h)--(m);\draw[->,thick] (m)--(d);\draw[->,thick] (d)--(v);\draw[->,thick] (v)--(r);
\end{tikzpicture}
\caption{Estructura mínima de una pregunta experimental para la trayectoria de investigación.}
\end{figure}

\section*{Cuatro distinciones que hacen investigable la pregunta}
\addcontentsline{toc}{section}{Del prototipo a la investigación aplicada}
\textbf{1. Variable observable.} La fracción de palabras partidas en subunidades
no es idéntica a los bits por carácter. La primera caracteriza la representación;
la segunda resume una pérdida probabilística si el modelo permite computarla.
\textbf{2. Tratamiento.} Entrenar un nuevo BPE modifica la segmentación; cambiar
simultáneamente arquitectura, número de parámetros y corpus no permite atribuir
una mejora solo a la segmentación.
\textbf{3. Unidad experimental.} Debe ser, cuando sea posible, el documento o
el informe, no cada subfragmento del mismo archivo; de lo contrario,
los intervalos basados en tokens individuales sobrestiman la precisión.
\textbf{4. Generalización.} Un vocabulario reducido para diez términos elegidos a mano
no es evidencia de efecto promedio en un dominio. Es una demostración inicial.

\begin{ejemplo}[Cómo pasar de observación a hipótesis falsable]
\textbf{Observación ilustrativa:} un tokenizador general parte el término
\emph{revestidor} en más fragmentos que otro entrenado en textos de petróleo.
\textbf{Hipótesis operacional futura:} sobre informes técnicos reservados,
con normalización y tamaño de vocabulario previamente fijados,
la diferencia media pareada de tokens por palabra es mayor que cero.
\textbf{Evaluación:} informar media, intervalo de confianza por remuestreo de
documentos y posibles fallas en términos que no estaban en el vocabulario.
No se informa un resultado numérico porque los datos no han sido descargados
ni el experimento ejecutado.
\end{ejemplo}

\section*{Fuentes candidatas, derechos y riesgos}
\addcontentsline{toc}{section}{Fuentes y límites de la evidencia}
\begin{center}\small
\begin{tabularx}{\linewidth}{@{}p{1.2cm}p{3.3cm}Xp{2.6cm}@{}}\toprule
ID & Fuente del catálogo & Uso inicial posible & Precaución\\\midrule
D01 & Wikipedia (eswiki) & Corpus amplio de artículos en español y temas científicos & Atribución y condiciones de reutilización según volcado/licencia\\
D02 & Proyecto Gutenberg & Obras en español para contraste con prosa técnica & Derecho de autor según país y obra; condiciones del portal\\
D03 & Volve, Equinor & Candidato posterior para texto y series de petróleo & Verificar presencia de texto, autorización y licencia de los archivos\\\bottomrule
\end{tabularx}
\end{center}
\textbf{Ubicaciones de origen:} \url{https://dumps.wikimedia.org/eswiki/},
\url{https://www.gutenberg.org/} y
\url{https://www.equinor.com/energy/volve-data-sharing}.
La fuente de la licencia debe verificarse antes de redistribuir; la existencia
de un catálogo no equivale a disponer de permiso para todo uso. La página de
condiciones de Gutenberg señala explícitamente la necesidad de examinar los
derechos en el país del usuario, además de los términos de la obra.

\section*{Conexión con Series de Tiempo y con otras áreas}
\addcontentsline{toc}{section}{Puente hacia Series de Tiempo}
El objeto probabilístico $p(x_{1:T})$ aparece en ambos cursos. Una cadena de
Markov utiliza una matriz de transición; un modelo de lenguaje tabular cuenta
símbolos y estima esas transiciones. La similitud no implica que un LLM sea una
cadena homogénea finita: sus condicionales pueden depender de una ventana amplia
y de una función parametrizada por una red.

Además de IC10, la matriz incluye problemas que conectan texto y pronóstico:
PG02 (eventos de perforación), FI01 (tono de noticias y volatilidad) y EE01
(demanda eléctrica con eventos textuales). Son \emph{posibilidades de exploración},
no compromisos de ejecución en la S01. En trabajos de pronóstico con modelos
preentrenados habrá que controlar también la fuga temporal y lo que el modelo
pueda haber memorizado; una partición cronológica de nuestros archivos no elimina
por sí sola ese riesgo.

\section*{Guion de discusión de la última hora}
\addcontentsline{toc}{section}{Preguntas que articulan la discusión}
\begin{enumerate}
\item ¿Qué hipótesis distingue una comparación de tokenizadores de una
 comparación de modelos completos? Dar una variable que se pueda medir.
\item ¿Qué evidencia refutaría la hipótesis en vez de confirmarla?
\item ¿Qué datos hay hoy, cuáles requieren autorización y qué licencias
 pueden condicionar la investigación?
\item ¿Qué cambia si la unidad de análisis es un documento y no un token?
\item ¿Cómo se formularía una pregunta que conecte texto fechado con una
 serie temporal, evitando que la información futura entre en el entrenamiento?
\end{enumerate}
\begin{investigacion}[Hito autónomo S01]
Cada estudiante entrega \textbf{dos áreas de interés}, un párrafo argumentado
por área. En cada párrafo se solicita: relevancia del problema, una posible
fuente de datos, una variable que eventualmente podría medirse y una cautela.
No se exige entrenar un modelo ni comprometer una hipótesis definitiva.
Esta secuencia anticipa los cinco candidatos de S02 y la decisión formal de S05.
\end{investigacion}

\section*{Del curso a una evidencia reproducible}
\addcontentsline{toc}{section}{Qué registrar desde la primera experiencia}
El repositorio del proyecto conservará configuración de entornos, fuentes,
transformaciones, semillas, pruebas y salidas regenerables. Las tablas de resultados
no deben editarse a mano. En ejercicios con IA generativa deberá declararse su
uso; toda derivación matemática deberá ser defendible y toda cita verificable.
Los proyectos con datos personales requieren procedimientos de anonimización y,
si corresponde, autorizaciones. Una herramienta que responde correctamente a una
pregunta de demostración no constituye una evidencia científica.
